% =========================================================
% Proof: Monotonicity of Supremum
% Source: volume-iii/analysis/bounding/notes/notes-bounding-monotonicty.tex
% =========================================================

\subsection*{Monotonicity of Supremum}
\label{prf:monotonicity-of-supremum}

\begin{remark*}[Return]
$\leftarrow$ \hyperref[thm:monotonicity]{Back to Theorem in Notes}
\end{remark*}

\begin{proof}
\Claim Let $A, B \subseteq \mathbb{R}$ be nonempty and bounded above.
If $A \subseteq B$, then $\sup A \le \sup B$.

\Given $A \subseteq B$; both sets bounded above.
\Goal $\sup A \le \sup B$.

\ByThm{Completeness Axiom} $\sup B$ exists.
\ByDef supremum, $b \le \sup B$ for all $b \in B$.

Since $A \subseteq B$, every $a \in A$ satisfies $a \in B$, hence $a \le \sup B$.
\Therefore $\sup B$ is an upper bound of $A$.

Since $\sup A$ is the least upper bound of $A$ and $\sup B$ is an upper bound
of $A$, we have $\sup A \le \sup B$. \AsReq
\end{proof}

\begin{remark*}[Proof shape]
One-step argument: identify $\sup B$ as an upper bound of $A$ (using
$A \subseteq B$), then invoke the least-upper-bound property of $\sup A$.
No $\varepsilon$-argument is needed.
\end{remark*}

\begin{remark*}[Dependencies]
Completeness Axiom (existence of $\sup B$); definition of supremum.
\end{remark*}
